% GENERATED FILE. Edit canonical lessons, then run npm run generate:books.
% canonical-source-hash: fnv1a64:5dbf629d3cd2b6ad
% canonical-lessons: FA-C125-seft, FA-C125-zesht, FA-C125-javan, FA-C125-pir, FA-C125-tanbal

\chapter{Hard, Ugly, Young, Old, Lazy}
\label{ch:fa-hard-ugly-young-old-lazy}
\hlchaptermodality{125}

\begin{chapteropening}
\textbf{By the end of this chapter:} \emph{I can say hard, ugly, young, old and lazy.}

Complete the last lesson of chapter 125: I can say hard, ugly, young, old and lazy.
\end{chapteropening}

\section[\texorpdfstring{\fa{سفت} (seft)}{سفت (seft)}]{\fa{سفت} (seft) — hard, firm}

\label{lesson:FA-C125-seft}



\begin{quote}
Before the new one: say the Persian for dry, then the Persian for soft.
\end{quote}

\subsection*{You'll want to know: \fa{سفت}}
\textbf{\fa{سفت}} — \emph{seft} — \textquotedblleft{}hard, firm\textquotedblright{}.

\begin{cousinweb}
\emph{seft} is a native Persian word, inherited rather than borrowed; it is usually traced to Middle Persian \emph{saft}, \textquotedblleft{}firm, hard\textquotedblright{}.
\end{cousinweb}

\begin{grammarlens}[title={What you've built}]
One more word for this chapter.
\end{grammarlens}

\subsection*{Your turn}
\begin{itemize}
\raggedright
  \item \emph{Say it:} \emph{seft}
  \item \emph{Say it:} \emph{seft}, once more
  \item \emph{Say it:} say \emph{seft}
  \item \emph{Recall:} say the Persian for dry, then the Persian for soft, then say \emph{seft} again
\end{itemize}

\begin{tcolorbox}[breakable,colback=teal!4,colframe=teal!35!black,title={Before you move on}]
What does \fa{سفت} mean? (\textquotedblleft{}Hard, firm\textquotedblright{}.) Say it once more.
\end{tcolorbox}
\section[\texorpdfstring{\fa{زشت} (zesht)}{زشت (zesht)}]{\fa{زشت} (zesht) — ugly}

\label{lesson:FA-C125-zesht}



\begin{quote}
Before the new one: say the Persian for soft, then the Persian for hard.
\end{quote}

\subsection*{You'll want to know: \fa{زشت}}
\textbf{\fa{زشت}} — \emph{zesht} — \textquotedblleft{}ugly\textquotedblright{}.

\begin{cousinweb}
\emph{zesht} is a native Persian word, inherited rather than borrowed: it continues Middle Persian \emph{zišt}, \textquotedblleft{}hateful, ugly\textquotedblright{} (EtymDB).
\end{cousinweb}

\begin{grammarlens}[title={What you've built}]
One more word for this chapter.
\end{grammarlens}

\subsection*{Your turn}
\begin{itemize}
\raggedright
  \item \emph{Say it:} \emph{zesht}
  \item \emph{Say it:} \emph{zesht}, once more
  \item \emph{Say it:} say \emph{zesht}
  \item \emph{Recall:} say the Persian for soft, then the Persian for hard, then say \emph{zesht} again
\end{itemize}

\begin{tcolorbox}[breakable,colback=teal!4,colframe=teal!35!black,title={Before you move on}]
What does \fa{زشت} mean? (\textquotedblleft{}Ugly\textquotedblright{}.) Say it once more.
\end{tcolorbox}
\section[\texorpdfstring{\fa{جوان} (javân)}{جوان (javân)}]{\fa{جوان} (javân) — young}

\label{lesson:FA-C125-javan}



\begin{quote}
Before the new one: say the Persian for hard, then the Persian for ugly.
\end{quote}

\subsection*{You'll want to know: \fa{جوان}}
\textbf{\fa{جوان}} — \emph{javân} — \textquotedblleft{}young\textquotedblright{}.

\begin{cousinweb}
\emph{javân} is a native Persian word, inherited rather than borrowed; it is usually traced to Middle Persian \emph{juwān}, \textquotedblleft{}young\textquotedblright{}. EtymDB takes it back to the Indo-European word for \textquotedblleft{}young\textquotedblright{} behind Latin \emph{iuvenis} and English \emph{young}.
\end{cousinweb}

\begin{grammarlens}[title={What you've built}]
One more word for this chapter.
\end{grammarlens}

\subsection*{Your turn}
\begin{itemize}
\raggedright
  \item \emph{Say it:} \emph{javân}
  \item \emph{Say it:} \emph{javân}, once more
  \item \emph{Say it:} say \emph{javân}
  \item \emph{Recall:} say the Persian for hard, then the Persian for ugly, then say \emph{javân} again
\end{itemize}

\begin{tcolorbox}[breakable,colback=teal!4,colframe=teal!35!black,title={Before you move on}]
What does \fa{جوان} mean? (\textquotedblleft{}Young\textquotedblright{}.) Say it once more.
\end{tcolorbox}
\section[\texorpdfstring{\fa{پیر} (pir)}{پیر (pir)}]{\fa{پیر} (pir) — old (of people)}

\label{lesson:FA-C125-pir}



\begin{quote}
Before the new one: say the Persian for ugly, then the Persian for young.

Say \textbf{\fa{اول}}, \textbf{\fa{دوم}}, \textbf{\fa{سوم}}, \textbf{\fa{چهارم}} in order. Then build \emph{panjom}, fifth, from \textbf{\fa{پنج}}.
\end{quote}

\subsection*{You'll want to know: \fa{پیر}}
\textbf{\fa{پیر}} — \emph{pir} — \textquotedblleft{}old (of people)\textquotedblright{}.

\begin{cousinweb}
\emph{pir} is a native Persian word, inherited rather than borrowed; it is usually traced to Middle Persian \emph{pīr}, \textquotedblleft{}old\textquotedblright{}. (EtymDB connects it with \emph{pedar}, \textquotedblleft{}father\textquotedblright{}, but that link is not settled.)
\end{cousinweb}

\begin{grammarlens}[title={What you've built}]
One more word for this chapter.
\end{grammarlens}

\subsection*{Your turn}
\begin{itemize}
\raggedright
  \item \emph{Say it:} \emph{pir}
  \item \emph{Say it:} \emph{pir}, once more
  \item \emph{Say it:} say \emph{pir}
  \item \emph{Recall:} say the Persian for ugly, then the Persian for young, then say \emph{pir} again
\end{itemize}

\begin{tcolorbox}[breakable,colback=teal!4,colframe=teal!35!black,title={Before you move on}]
What does \fa{پیر} mean? (\textquotedblleft{}Old (of people)\textquotedblright{}.) Say it once more.
\end{tcolorbox}
\section[\texorpdfstring{\fa{تنبل} (tanbal)}{تنبل (tanbal)}]{\fa{تنبل} (tanbal) — lazy}

\label{lesson:FA-C125-tanbal}



\begin{quote}
Before the new one: say the Persian for young, then the Persian for old.
\end{quote}

\subsection*{You'll want to know: \fa{تنبل}}
\textbf{\fa{تنبل}} — \emph{tanbal} — \textquotedblleft{}lazy\textquotedblright{}.

\begin{cousinweb}
\emph{tanbal} is a native Persian word (Turkish \emph{tembel} comes from it), but its older origin is not settled.
\end{cousinweb}

\begin{grammarlens}[title={What you've built}]
One more word for this chapter.
\end{grammarlens}

\subsection*{Your turn}
\begin{itemize}
\raggedright
  \item \emph{Say it:} \emph{tanbal}
  \item \emph{Say it:} \emph{tanbal}, once more
  \item \emph{Say it:} say \emph{tanbal}
  \item \emph{Recall:} say the Persian for young, then the Persian for old, then say \emph{tanbal} again
\end{itemize}

\begin{tcolorbox}[breakable,colback=teal!4,colframe=teal!35!black,title={Before you move on}]
What does \fa{تنبل} mean? (\textquotedblleft{}Lazy\textquotedblright{}.) Say it once more.
\end{tcolorbox}
